\addsec{Conclusion}
\markboth{Conclusion}{}

Le système obtenu associe un processeur Verilog à quatre RAM VHDL de 1024 mots de 32 bits.
Le même modèle de RAM est utilisé seul, dans les tests de \texttt{quadram} et dans le
système final. Les paramètres permettent d'en changer les dimensions sans réécrire son
architecture.

Les essais ont surtout mis en évidence deux conditions d'intégration : seule la RAM
sélectionnée doit accepter une écriture, et le processeur doit attendre la mémorisation
de l'adresse avant de lire une donnée. Le programme \texttt{test\_memoire} vérifie la
sélection des quatre blocs, un déplacement négatif et une frontière entre deux RAM.
La comparaison de \texttt{instr.txt} et \texttt{instr2.txt} montre aussi que déplacer les
données vers une mémoire externe oblige à revoir leur initialisation.

Les tests de validation du processeur corrigé et du système donnent les résultats attendus.
La vérification finale reproduit également les défauts du modèle d'origine et l'échec du
premier programme lorsque les données externes ne sont pas initialisées. Cette validation
reste celle des cas testés : le processeur implémente un sous-ensemble de RV32I et les huit
instructions de l'extension M, et non toute l'ISA RV32IM.

Pour aller jusqu'à une implantation sur FPGA, il faudrait remplacer les attentes
comportementales par une machine à états synchrone, utiliser une sélection de données
adaptée aux ressources internes de la cible, puis vérifier la synthèse, les contraintes
d'horloge et le fonctionnement sur carte. Pour exécuter du code issu d'un compilateur C
RISC-V, il faudrait compléter les instructions, passer à l'adressage par octets et prévoir
le démarrage du programme, la pile et la répartition des mémoires. Ces deux prolongements
n'ont pas été réalisés.
